\documentclass[10pt,a4paper]{amsart}
\usepackage[margin=19mm]{geometry}
\usepackage{amsmath,amssymb,mathtools}
\usepackage[scr=rsfs,cal=euler]{mathalfa}
\usepackage{xfrac}
\usepackage[unicode,hidelinks,bookmarks=false]{hyperref}
\newcommand{\ie}{i.e.\ }
\newcommand{\cf}{cf.\ }
\newcommand{\resp}{respectively\ }
\newcommand{\Subsection}{\subsection}
\providecommand{\nobreakdash}{\nobreak\hbox{-}}
\setlength{\emergencystretch}{2em}
\multlinegap0pt
\newcommand{\cal}{\mathcal}
\newtheorem{theorem}{Theorem}[section]
\newtheorem{proposition}[theorem]{Proposition}
\newtheorem{conjecture}[theorem]{Conjecture}
\newtheorem{lemma}[theorem]{Lemma}
\newtheorem{corollary}[theorem]{Corollary}
\newtheorem{convention}{Convention}[theorem]
\theoremstyle{definition}
\newtheorem{def-theorem}[theorem]{Definition-Theorem}
\newtheorem{remark}[theorem]{Remark}
\newtheorem{definition}[theorem]{Definition}
\newtheorem{example}[theorem]{Example}
\title[Finite entropy sums in quantum field theory]{Finite entropy sums in quantum field theory: the Fourier-coefficient criterion for safe functions on a hypergraph (Lemma 4.2)}
\author{Mark Van Raamsdonk}
\address{Department of Physics and Astronomy, University of British Columbia,\\ 6224 Agricultural Road, Vancouver, B.C., V6T 1Z1, Canada}
\date{Source: arXiv:2508.21276v1 (CC BY 4.0)}
\begin{document}
\begin{abstract}
Entropies associated with spatial subsystems in conventional local quantum field theories are typically divergent when the spatial regions have boundaries. However, in certain linear combinations of the entropies for various subsystems, these divergences may cancel, giving finite quantities that provide information-theoretic data about the underlying state. In this note, we show that all such quantities can be written as linear combinations of three basic types of quantities: i) the entropy of a spatial subsystem minus the entropy of its complementary subsystem, ii) the mutual information between non-adjacent subsystems, and iii) the tripartite information for triples of disjoint subsystems. For a fixed decomposition of a spatial slice into regions, we describe a basis of sums of entropies for collections of for these regions for which all divergences related to both region boundaries and higher-codimension intersections of regions cancel. Key mathematical technology used in this work (Fourier transforms on the Boolean cube and M\"obius transformations of functions on partially ordered sets) and several of the main proof ideas were suggested by AI (ChatGPT5). We offer a few comments on the use of AI in physics and mathematics, based on our experience. 
\end{abstract}
\maketitle

\noindent\emph{Source and licence.} Finite entropy sums in quantum field theory. Version arXiv:2508.21276v1 (CC BY 4.0), distributed under the Creative Commons Attribution 4.0 International licence, http://creativecommons.org/licenses/by/4.0/, as recorded in the arXiv OAI licence metadata for this exact version and shown on the abstract page https://arxiv.org/abs/2508.21276v1. Full rights notice: licenses/arxiv-2508.21276-CC-BY-4.0.txt.\par\smallskip
\noindent\textit{Editorial scope.} Counted unit: the paper's Fourier-coefficient criterion for the space of safe functions on a hypergraph, printed as Lemma 4.2 (source0.tex lines 557--560), delivered together with the paper's complete original proof of it at source0.tex lines 571--621 as the single primary proof body. No mathematics is written, repaired or re-derived: statement and proof are the paper's own text line for line, in the paper's own notation (the script T for the space of real functions on subsets of the vertices, its subspace of safe functions, the intersection family, the Fourier basis and its coefficients, the test function of a bipartition, and the parity conventions of the paper). The counted proof environment carries the paper's own annotation naming the lemma it proves, and it is the only proof of the counted lemma in the paper; the theorem printed immediately after the lemma has no proof of its own and is not part of this unit. Required context retained uncounted, in document order: the opening of Section 4 with the motivation for higher-codimension intersections, the generalisation from a graph to a hypergraph, the definition of the family of intersections and the definition of the subspace of safe functions with its numbered condition (13) and the equivalent test-function formulation (source0.tex lines 530--556); and the lead-in sentence together with the statement of the theorem that the counted lemma immediately yields, so that the reader sees what the lemma is for (source0.tex lines 561--570). Every definition, numbered display, footnote and citation of those lines is retained. Not part of this unit and not redistributed: the paper's front matter (source0.tex lines 59--82, of which the title, the author name and the affiliation are reproduced in the wrapper title block and the abstract is reproduced in the wrapper front matter); the abstract as printed in the source (source0.tex line 84, reproduced as discussed below); the whole of Sections 1 to 3 and of Sections 5 and 6, that is the introduction, the graph-theory formulation, the codimension-one cancellation analysis with its Fourier, information and Mobius bases, the quantum-mechanics reformulation and the algebraic definitions (source0.tex lines 96--529 and 622--780), including the proofs of the lemmas and propositions on which the paper's earlier results rest; the acknowledgements; and the paper's bibliography listing, of which only the cited entries are transcribed. The delivered text cites thirteen keys through exactly two citation commands, the twelve-key citation of the references for higher-codimension divergences in the opening paragraph of Section 4 and the single-key citation in the footnote attached to the same paragraph, and those thirteen entries are transcribed verbatim from the external bibliography file that the paper's own build ships (the file main.bbl in the e-print, thirty entries in all), with the printed labels preserved as 8 to 20, so that the delivered citation prints the same reference numbers as the published version, namely the range 8 to 19 for the twelve-key citation and 20 for the footnote citation. One presentational difference is declared: the wrapper does not load the citation-formatting package that the paper loads, so that twelve-key citation prints as an explicit list of the numbers 8 to 19 rather than collapsed into the range 8--19. The numbers themselves, and therefore every reference the reader is sent to, are identical. Reference adaptation: exactly one was needed and exactly one was made. The retained lead-in sentence that announces the theorem following the counted lemma refers to the paper's earlier graph-theoretic Fourier theorem, whose label lies outside every retained span; that single cross-reference is delivered as a hyperlink to the pinned version with the published number kept literally, so the delivered sentence reads 'following the same logic as in the proof of Theorem 3.3' with the number 3.3 hyperlinked. All other labels and references inside the delivered unit resolve inside it: the label of the numbered condition (13) of the retained definition, which the counted proof cites when it sets up its Fourier expansion, and the three equation labels introduced inside the counted proof itself, which the proof cites in its own final steps. The paper's own counters are reproduced so that the printed numbers coincide with the published version: Section 4 and its first subsection, Definition 4.1 for the retained definition of the safe subspace, Lemma 4.2 for the counted lemma, Theorem 4.3 for the theorem printed after it, and equations (13) to (16) for the retained and the proof-internal numbered displays, which is why the wrapper numbers equations sequentially from the paper's own starting value rather than within sections. Printed slips kept literal: no printed word, symbol, hypothesis or number of the counted statement or of its proof was altered. The paper's own forms are delivered as printed, among them its use of the old-style calligraphic alphabet for the space of functions, its printed misspelling of the word product in one sentence of the counted proof, its notation for the intersection family and for the even-order subsets of an intersection, its habit of writing the two-sided parity condition as a single equation, and its own reference keys. Non-ASCII characters: the source contains characters outside ASCII in one distinct code point only, the o-umlaut of Mobius, which occurs in the abstract and in a handful of body lines that lie outside every retained span. The retained spans contain no character outside ASCII at all. In the abstract, which is reproduced in the wrapper front matter, that o-umlaut is delivered with the standard LaTeX accent command rather than as a literal character, because the wrapper keeps the class's own text encoding, in which the literal character has no glyph under XeTeX; the printed result is the same letter and no character is dropped, replaced or re-encoded, and the compile receipt reports no missing character. Layout and class: the paper's own class is article at 12pt with a titlepage block, and it loads graphicx, amsmath, amsthm, amsfonts, courier, amssymb and a citation-formatting package; the wrapper is the lane's pinned amsart editorial wrapper at 10pt on A4, which already supplies amsmath, amssymb and mathtools, to which this package adds the paper's own theorem declarations and its definition-style switch copied verbatim. The paper's loads of graphicx, courier and the citation package are not reproduced, and neither are the old-style equation and array shorthand macros of its preamble, because no retained span uses them, with one exception: the paper writes its calligraphic letters for the space of functions and for the family of intersections with the old-style calligraphic font switch, which the wrapper's pinned mathematics-alphabet package does not provide, so this package defines that switch to expand to the calligraphic alphabet; the wrapper has pinned that alphabet to Euler script where the paper's own preamble used the default calligraphic alphabet, a purely typographic substitution that changes no mathematics. The two letters affected are the script I of the family of intersections and the script T of the space of real functions on subsets of the vertices. That is a page-layout and citation-format difference only, and no formula, symbol, hypothesis or argument changes. No publisher class, no publisher style file and no upstream script is redistributed. Assets: the paper has five figure environments whose five images are external image files, and it also names two of those files with a letter case that does not match the names actually shipped in the e-print archive; none of the five figure environments and none of the five image names lies inside any retained span, no retained span contains a graphics-inclusion command, an input directive or an include directive, the package declares no asset dependency, no image file is copied into this package, and no image file is redistributed with it. All mathematics of the unit is therefore native text with no image dependency. A rights notice is delivered at licenses/arxiv-2508.21276-CC-BY-4.0.txt.
\setcounter{section}{3}\setcounter{subsection}{0}\setcounter{theorem}{0}\renewcommand{\theequation}{\arabic{equation}}
\section{Cancelling divergences from multiple intersections}

\setcounter{equation}{12}

So far, we have been concerned about cancellation of divergences related to the pairwise intersections of regions. Our considerations apply equally well to the leading area-law divergences and to any subleading divergences or divergences associated with non-smooth features of these boundary regions, for example corners in a one dimensional boundary between a pair of regions in two dimensions. To see this, we can use the fact that the divergences appearing in entropies are independent of the state (they arise from the UV degrees of freedom in the quantum field theory whose entanglement structure is the same in different finite-energy states). For a pure state of the whole system, the entropy of a region $A$ equals the entropy of the complementary region $A^c$, so in particular, all divergences associated with the boundary of $A$ must match the divergences associated with the boundary of $A^c$. Thus, our condition that the sum of coefficients for all regions having a particular boundary component (regardless of which side of the boundary the region occupies) should ensure that all types of divergences associated with this boundary component will cancel. 

Additional divergences may arise from the intersection of $m>2$ regions at some higher codimension feature $J$ \cite{CasiniHuerta2007_Universal2p1,CasiniHuerta2009_Review,BuenoMyersWitczak2015_PRL,BuenoMyers2015_JHEP,BuenoMyersWitczak2015_Twist,KallinEtAl2014_JStatMech,StoudenmireEtAl2014_PRB,HaywardSierensEtAl2017_TrihedralPRB,BednikEtAl2019_TrihedralPRB,SeminaraSistiTonni2017_BCFT,Berthiere2019_BoundaryCornerPRB,Nishioka2018_RMP}. Let $I_m$ be the subset of vertices corresponding to the regions $A_i$ whose boundaries intersect at $J$. For any non-empty proper subset $\tau \subset I_m$, the region $A_\tau$ has a potential divergence associated with the ``corner'' of $A_\tau$ that touches $J$.\footnote{It is possible that no such divergence appears: for example, with three two-dimensional regions intersecting at a point, one of the three angles might be $\pi$ in which case there is no corner divergence for the corresponding region. It may also be that the corner divergence is absent if the corner is sufficiently ``cuspy'' \cite{BuenoCasiniWitczak2019_SingularJHEP}.} In an entropy sum, this divergence (if it exists) will be cancelled provided that the sum of coefficients for sets $\sigma$ containing $\tau$ but not $I_m \backslash \tau$ (the complement of $\tau$ in $I_m$) plus the sum of coefficients for sets $\sigma$ containing $I_m \backslash \tau$ but not $\tau$ vanishes. Here, we are using the fact that the divergence associated with a corner will be the same as the divergence associated with its complementary corner. As above, this follows since divergences are state independent, and the entropy of a subsystem is the same as the entropy of its complement in a pure state. 

We can generalize our earlier analysis in order to analyze these additional conditions associated with multiple intersections. In order to capture the information about all intersections, we can generalize from a graph to a {\it hypergraph} $H$, which is again a set $H_1 \equiv V$ of $n$ vertices associated with the $n$ regions, but now with sets $H_2 \equiv E, H_3, H_4 \dots$, of pairs, triples, quadruples, etc... of vertices corresponding to the intersections of 3,4, etc... regions. We define ${\cal I}_H = \cup_{m \ge 2} H_m$ and refer to this as the set of all {\it intersections}.

Starting from the space ${\cal T}_\mathbb{R}$ of functions on subsets of the vertices, we can now define a subspace ${\cal T}_H$ of functions that are safe according to the conditions above:
\begin{definition}
    Let $H$ be a hypergraph and ${\cal T}_{\mathbb{R}}$ be the space of real functions on subsets of its vertices. We define the subspace ${\cal T}_H$ of {\it safe} functions by the condition that for each intersection $I \in {\cal I}_H$ and each bipartition $\{a,b\}$ of $I$, 
\begin{equation}
\label{gencond}
\sum_{\sigma | a \in \sigma, b \notin \sigma} T(\sigma) + \sum_{\sigma | a \notin \sigma, b \in \sigma} T(\sigma) = 0 \; .
\end{equation}
\end{definition}
Equivalently, defining functions
\[
g_{a,b}(\sigma) = \left\{ \begin{array}{ll} 1 & (a \cup b) \cap \sigma \in \{a,b\} \cr
0 & \qquad {\rm otherwise} \end{array}  \right.
\]
for nonintersecting sets $a,b$, ${\cal T}_H$ is the subspace defined by $(g_{a,b}, T) = 0$ for all bipartitions $\{a,b\}$ of all intersections $I$. 

\subsection{A Fourier basis for ${\cal T}_H$}

We will now see that, as for the edge constraints, the constraints related to general intersections may be expressed very simply in the Fourier basis:


\begin{lemma}
\label{FourLemma}
For a general hypergraph $H = \{H_m\}$, a function $T \in {\cal T}_\mathbb{R}$ is in ${\cal T}_H$ if and only if for any intersection $I \in {\cal I}_H$ and any even order subset $\tau \subset I$, the Fourier coefficient $\hat{T}(\tau)$ is equal to $\hat{T}(\emptyset)$. 
\end{lemma}

This immediately leads (following the same logic as in the proof of Theorem \href{https://arxiv.org/abs/2508.21276v1}{3.3}) to
\begin{theorem}\label{thm:Fourier2}
For a hypergraph $H$, define $I_E$ to be the collection of even order subsets of vertices (including the empty set) that are subsets of some intersection $I \in {\cal I}_H$. Then an orthogonal basis for ${\cal T}_H$ is provided by the Fourier basis elements $\chi_\sigma$ for $\sigma \notin I_E$ together with the function
\[
\chi_0 \equiv \sum_{\sigma \in I_E} \chi_\sigma \; . 
\]
The dimension of ${\cal T}_H$ is $2^n - |I_E| + 1$.
\end{theorem}



\begin{proof} (Lemma \ref{FourLemma})
Consider an intersection $I$ and some bipartition $\{a,b\}$ of $I$.
Expressed in terms of the Fourier basis, the condition (\ref{gencond}) for the pair $\{a,b\}$ is
\begin{equation}
\label{multicond}
\sum_\sigma \hat{T}(\sigma) (g_{a,b}, \chi_\sigma) = 0 \; ,
\end{equation}
where
\[
(g_{a,b}, \chi_\sigma) = \sum_\tau g_{a,b}(\tau)  (-1)^{x_\sigma \cdot x_\tau} \; .
\]
Now, for any vertex $i$ not in $I$ and any set $\tau$ for which $g_{a,b} \ne 0$, there is a distinct complementary set $\tau'$ with $g_{a,b} \ne 0$ obtained by removing / adding the element $i$ if it is present / not present. If $i$ is in $\sigma$, then the terms in the sum above corresponding to complementary sets $\tau$ and $\tau'$ will cancel each other, so the sum vanishes. Thus, inner prduct above is nonzero only for $\sigma \subset I$. 

The terms in the sum for which $g_{a,b}(\tau)$ is nonvanishing are of the form $\tau = a \cup \varphi$ or $\tau = b \cup \varphi$ for $\varphi \cap I = \emptyset$. So we have, for $\sigma \subset I$,
\[
(g_{a,b}, \chi_\sigma) = \sum_\varphi (-1)^{x_\sigma \cdot x_{a \cup \varphi}} + (-1)^{x_\sigma \cdot x_{b \cup \varphi}} .
\]
Since $\sigma \cap \varphi = \emptyset$, we have $x_\sigma \cdot x_{a \cup \varphi} = x_\sigma \cdot x_{a}$, and $x_\sigma \cdot x_{b \cup \varphi} = x_\sigma \cdot x_{b}$. Thus, for any $\sigma$,
\[
(g_{a,b}, \chi_\sigma) =  \left\{ \begin{array}{ll} 2^{n-|I|} ((-1)^{x_\sigma \cdot x_a} + (-1)^{x_\sigma \cdot x_b}) &  \qquad \sigma \subset I \cr
0 & \qquad {\rm otherwise} \end{array}  \right.
\]
It follows that the condition (\ref{multicond}) involves only Fourier coefficients $\hat{T}_\sigma$ for sets $\sigma \subset I$, reducing to 
\begin{equation}
\label{multicond2}
\sum_{\sigma \subset I} \hat{T}(\sigma) ((-1)^{x_\sigma \cdot x_a} + (-1)^{x_\sigma \cdot x_b}) = 0 \; .
\end{equation}
The expression in brackets is nonzero if and only if $a \cap \sigma$ and $b \cap \sigma$ have the same parity, which is equivalent to $|\sigma| = |(a \cap \sigma) \cup (b \cap \sigma)|$ being even. We can now further reduce the condition to
\[
\sum_{\sigma \subset I, |\sigma| \; even}  \hat{T}(\sigma) (-1)^{x_\sigma \cdot x_a} = 0
\]
Since we have one equation for each bipartition of $I$, this is a set of $2^{|\sigma|-1}-1$ equations for $2^{|\sigma|-1}$ unknowns. We will show that these equations are all independent, leaving a one-dimensional space of solutions in which all $\hat{T}$s appearing here are equal. 

Choose some arbitrary element $1 \in I$ and let $\tilde{I} = I \backslash 1$. Then bipartitions $\{a,b\}$ are in one-to-one correspondence with non-empty subsets $c \in \tilde{I}$ via $a= c, b = I \backslash c$. Define a bijection $e$ from subsets of $\tilde{I}$ to even order subsets of $I$ by 
\[
e(\sigma) = \left\{ \begin{array}{ll} \sigma & \qquad |\sigma| \; even \cr
{1} \cup \sigma & \qquad |\sigma| \; odd \cr
\end{array}
\right.
\]
Then our condition is equivalent to the statement that for every $c \in \tilde{I}$
\begin{equation}
\label{reduced}
\sum_{\sigma \subset \tilde{I}} \hat{T}(e(\sigma)) (-1)^{y_\sigma \cdot y_c} = 0
\end{equation}
where $y_\sigma \in \mathbb{R}^{|I| - 1}$ is the characteristic function for a subset $\sigma \in \tilde{I}$. Defining $\{\tilde{\chi}_\sigma\}$ to be the Fourier basis for functions on subsets of $\tilde{I}$ and defining a function $C$ on subsets of $\tilde{I}$ by
\[
C = \sum_\sigma \hat{T}(e(\sigma)) \tilde{\chi}_\sigma \; ,
\]
the condition (\ref{reduced}) is simply that $C(c) = 0$ for all non-empty $c \subset \tilde{I}$. This is equivalent to the equality of all the Fourier coefficients $\hat{T}(e(\sigma))$, which are the Fourier coefficients associated with the even subsets of $I$. 
\end{proof}

\begin{thebibliography}{99}
\bibitem[8]{CasiniHuerta2007_Universal2p1} H.~Casini and M.~Huerta, \emph{Universal terms for the entanglement entropy in 2+1 dimensions}, \href{http://dx.doi.org/10.1016/j.nuclphysb.2006.12.012}{\emph{Nuclear Physics B} {\bf 764} (2007) 183--201}, [\href{https://arxiv.org/abs/hep-th/0606256}{{\tt hep-th/0606256}}].
\bibitem[9]{CasiniHuerta2009_Review} H.~Casini and M.~Huerta, \emph{Entanglement entropy in free quantum field theory}, \href{http://dx.doi.org/10.1088/1751-8113/42/50/504007}{\emph{Journal of Physics A: Mathematical and Theoretical} {\bf 42} (2009) 504007}, [\href{https://arxiv.org/abs/0905.2562}{{\tt 0905.2562}}].
\bibitem[10]{BuenoMyersWitczak2015_PRL} P.~Bueno, R.~C. Myers and W.~Witczak-Krempa, \emph{Universality of corner entanglement in conformal field theories}, \href{http://dx.doi.org/10.1103/PhysRevLett.115.021602}{\emph{Physical Review Letters} {\bf 115} (2015) 021602}, [\href{https://arxiv.org/abs/1505.04804}{{\tt 1505.04804}}].
\bibitem[11]{BuenoMyers2015_JHEP} P.~Bueno and R.~C. Myers, \emph{Corner contributions to holographic entanglement entropy}, \href{http://dx.doi.org/10.1007/JHEP08(2015)068}{\emph{Journal of High Energy Physics} {\bf 2015} (2015) 68}, [\href{https://arxiv.org/abs/1505.07842}{{\tt 1505.07842}}].
\bibitem[12]{BuenoMyersWitczak2015_Twist} P.~Bueno, R.~C. Myers and W.~Witczak-Krempa, \emph{Universal corner entanglement from twist operators}, \href{http://dx.doi.org/10.1007/JHEP09(2015)091}{\emph{Journal of High Energy Physics} {\bf 2015} (2015) 91}, [\href{https://arxiv.org/abs/1507.06997}{{\tt 1507.06997}}].
\bibitem[13]{KallinEtAl2014_JStatMech} A.~B. Kallin, E.~M. Stoudenmire, P.~Fendley, R.~R.~P. Singh and R.~G. Melko, \emph{Corner contribution to the entanglement entropy of an o(3) quantum critical point in 2+1 dimensions}, \href{http://dx.doi.org/10.1088/1742-5468/2014/06/P06009}{\emph{Journal of Statistical Mechanics: Theory and Experiment} (2014) P06009}, [\href{https://arxiv.org/abs/1401.3504}{{\tt 1401.3504}}].
\bibitem[14]{StoudenmireEtAl2014_PRB} E.~M. Stoudenmire, P.~Gustainis, R.~Johal, S.~Wessel and R.~G. Melko, \emph{Corner contribution to the entanglement entropy of strongly interacting o(2) quantum critical systems in 2+1 dimensions}, \href{http://dx.doi.org/10.1103/PhysRevB.90.235106}{\emph{Physical Review B} {\bf 90} (2014) 235106}, [\href{https://arxiv.org/abs/1409.6327}{{\tt 1409.6327}}].
\bibitem[15]{HaywardSierensEtAl2017_TrihedralPRB} L.~E.~H. Sierens, P.~Bueno, R.~R.~P. Singh, R.~C. Myers and R.~G. Melko, \emph{Cubic trihedral corner entanglement for a free scalar}, \href{http://dx.doi.org/10.1103/PhysRevB.96.035117}{\emph{Physical Review B} {\bf 96} (2017) 035117}, [\href{https://arxiv.org/abs/1703.03413}{{\tt 1703.03413}}].
\bibitem[16]{BednikEtAl2019_TrihedralPRB} G.~Bednik, L.~E.~H. Sierens, M.~Guo, R.~C. Myers and R.~G. Melko, \emph{Probing trihedral corner entanglement for dirac fermions}, \href{http://dx.doi.org/10.1103/PhysRevB.99.155153}{\emph{Physical Review B} {\bf 99} (2019) 155153}, [\href{https://arxiv.org/abs/1810.02831}{{\tt 1810.02831}}].
\bibitem[17]{SeminaraSistiTonni2017_BCFT} D.~Seminara, J.~Sisti and E.~Tonni, \emph{Corner contributions to holographic entanglement entropy in ads$_4$/bcft$_3$}, \href{http://dx.doi.org/10.1007/JHEP11(2017)076}{\emph{Journal of High Energy Physics} {\bf 2017} (2017) 76}, [\href{https://arxiv.org/abs/1708.05080}{{\tt 1708.05080}}].
\bibitem[18]{Berthiere2019_BoundaryCornerPRB} C.~Berthiere, \emph{Boundary-corner entanglement for free bosons}, \href{http://dx.doi.org/10.1103/PhysRevB.99.165113}{\emph{Physical Review B} {\bf 99} (2019) 165113}, [\href{https://arxiv.org/abs/1811.12875}{{\tt 1811.12875}}].
\bibitem[19]{Nishioka2018_RMP} T.~Nishioka, \emph{Entanglement entropy: holography and renormalization group}, \href{http://dx.doi.org/10.1103/RevModPhys.90.035007}{\emph{Reviews of Modern Physics} {\bf 90} (2018) 035007}, [\href{https://arxiv.org/abs/1801.10352}{{\tt 1801.10352}}].
\bibitem[20]{BuenoCasiniWitczak2019_SingularJHEP} P.~Bueno, H.~Casini and W.~Witczak-Krempa, \emph{Generalizing the entanglement entropy of singular regions in conformal field theories}, \href{http://dx.doi.org/10.1007/JHEP08(2019)069}{\emph{Journal of High Energy Physics} {\bf 2019} (2019) 69}, [\href{https://arxiv.org/abs/1904.11495}{{\tt 1904.11495}}].
\end{thebibliography}
\end{document}
